\section{Space Complexity --- Exp BFS --- Run}
%%% generated by the blueprint pipeline from TCSlib.Complexity.SpaceComplexity.ExpBFS.Run
\subsection*{Space Complexity --- Exp BFS --- Run: Introduction}

Declarations of the Lean module \texttt{TCSlib.Complexity.SpaceComplexity.ExpBFS.Run}.

\subsection*{Space Complexity --- Exp BFS --- Run: Results}

\begin{definition}[Deterministic Advice Program]
\lean{Complexity.LogProg.ExpBFS.Ad}
\label{Complexity.LogProg.ExpBFS.Ad}
\leanok
\uses{Complexity.LogProg.ARM, Complexity.LogProg.ExpBFS.A, Complexity.LogProg.ExpBFS.Lb, Complexity.LogProg.MARM.detPart}
For each natural number $K_0$, let $A(K_0)$ be the program family specified by the
surrounding construction. The abbreviation denotes its deterministic part with the
Boolean parameter set to false, viewed as a register program with 17 registers, 4 oracle
deciders, and labels from $Lb(K_0)$.
\end{definition}

\begin{definition}[Machine Configuration]
\lean{Complexity.LogProg.ExpBFS.cf}
\label{Complexity.LogProg.ExpBFS.cf}
\leanok
\uses{Complexity.LogProg.ExpBFS.Lb, Complexity.LogProg.ExpBFS.St, Complexity.LogProg.ExpBFS.St.toV, Complexity.LogProg.MConf}
Given a label $l\in Lb(K_0)$, a state $st$, a memory function $M:\mathbb
Z\to\{\mathrm{false},\mathrm{true}\}$, and a memory pointer $p\in\mathbb Z$, the
abbreviation denotes the machine configuration whose active instruction is $l$, whose
register values are those encoded by $st$, whose answer is unset, and whose memory
consists of $M$ with pointer $p$.
\end{definition}

\begin{lemma}[Radius register differs from the source register]
\lean{Complexity.LogProg.ExpBFS.rad_ne_U}
\label{Complexity.LogProg.ExpBFS.rad_ne_U}
\leanok
\uses{Complexity.LogProg.ExpBFS.Ph3, Complexity.LogProg.ExpBFS.rU, Complexity.LogProg.ExpBFS.rad}
\difficulty{4}
For every phase $p$ among the first search, the second search, and the final scan, the
source register is different from the register assigned as the radius in phase $p$.
\end{lemma}

\begin{lemma}[Radius register differs from the target register]
\lean{Complexity.LogProg.ExpBFS.rad_ne_V}
\label{Complexity.LogProg.ExpBFS.rad_ne_V}
\leanok
\uses{Complexity.LogProg.ExpBFS.Ph3, Complexity.LogProg.ExpBFS.rV, Complexity.LogProg.ExpBFS.rad}
\difficulty{4}
For every phase $p$ among the first search, the second search, and the final scan, the
target register is different from the register assigned as the radius in phase $p$.
\end{lemma}

\begin{lemma}[Validity of the search-machine precondition]
\lean{Complexity.LogProg.ExpBFS.preS_A}
\label{Complexity.LogProg.ExpBFS.preS_A}
\leanok
\uses{Complexity.LogProg.AConf, Complexity.LogProg.ExpBFS.A, Complexity.LogProg.ExpBFS.Ad, Complexity.LogProg.ExpBFS.Lb, Complexity.LogProg.ExpBFS.Site.set, Complexity.LogProg.ExpBFS.acI, Complexity.LogProg.ExpBFS.cpI, Complexity.LogProg.ExpBFS.mvI, Complexity.LogProg.ExpBFS.rad, Complexity.LogProg.ExpBFS.rad_ne_U, Complexity.LogProg.ExpBFS.shI, Complexity.LogProg.ExpBFS.toPh, Complexity.LogProg.MARM.detPart, Complexity.LogProg.PreS}
\difficulty{6}
For every Boolean string $x$ and every abstract configuration of the search machine with
parameter $K_0$, the machine’s precondition holds: any register comparison uses distinct
registers, and any oracle call uses a list of pairwise distinct registers.
\end{lemma}

\begin{definition}[Guarded Configuration Predicate]
\lean{Complexity.LogProg.ExpBFS.G}
\label{Complexity.LogProg.ExpBFS.G}
\leanok
\uses{Complexity.LogProg.ExpBFS.Ad, Complexity.LogProg.ExpBFS.Lb, Complexity.LogProg.GuardLe, Complexity.LogProg.MConf}
For a word $x\in\{0,1\}^*$, natural numbers $B,P$, and machine configuration $c$,
$G(x,B,P,c)$ holds exactly when the first component of $c$ satisfies the guard predicate
for the program $Ad(K_0)$ on input $x$ with bound $B$, and the pointer in the memory
component of $c$ lies between $0$ and $P$, inclusive.
\end{definition}

\begin{lemma}[Bounded Initial Configuration]
\lean{Complexity.LogProg.ExpBFS.g_of}
\label{Complexity.LogProg.ExpBFS.g_of}
\leanok
\uses{Complexity.LogProg.ExpBFS.G, Complexity.LogProg.ExpBFS.Lb, Complexity.LogProg.ExpBFS.St, Complexity.LogProg.ExpBFS.St.Bd, Complexity.LogProg.ExpBFS.St.Bd.toV, Complexity.LogProg.ExpBFS.cf, Complexity.LogProg.ExpBFS.preS_A}
\difficulty{1}
Let $K_0,B,P$ be natural numbers, let $x$ be a Boolean word, let $\text{orc}$ assign a
Boolean answer to each of four indices and each Boolean word, and let
$M\colon\mathbb{Z}\to\{\mathrm{false},\mathrm{true}\}$ be a memory marking. Suppose a
machine state $st$ has all 17 register values at most $B$, and an integer $p$ satisfies
$0\le p\le P$. Then, for every program label $l$, the configuration with label $l$,
registers from $st$, no answer yet, memory $M$, and memory head at $p$ satisfies the
invariant: its instruction safety condition holds, all register values are at most $B$,
and the head lies between $0$ and $P$, inclusive.
\end{lemma}

\begin{lemma}[Deterministic Step Preserves Reachability]
\lean{Complexity.LogProg.ExpBFS.step_of}
\label{Complexity.LogProg.ExpBFS.step_of}
\leanok
\uses{Complexity.LogProg.ExpBFS.A, Complexity.LogProg.ExpBFS.G, Complexity.LogProg.ExpBFS.Lb, Complexity.LogProg.ExpBFS.St, Complexity.LogProg.ExpBFS.St.Bd, Complexity.LogProg.ExpBFS.St.toV, Complexity.LogProg.ExpBFS.cf, Complexity.LogProg.ExpBFS.g_of, Complexity.LogProg.Ins, Complexity.LogProg.MConf, Complexity.LogProg.MReach, Complexity.LogProg.MReach.step, Complexity.LogProg.astep, Complexity.LogProg.mstep}
\difficulty{3}
Let $K_0,B,P\in\mathbb{N}$, let $\mathcal O$ assign a Boolean oracle on finite bit
strings to each of four indices, and let $x$ be a finite bit string. Let $l$ be a
machine label with deterministic instruction $i$, let $st$ be a state whose 17 register
values are all at most $B$, let $M\colon\mathbb{Z}\to\{0,1\}$ be a memory tape, and let
$p\in\mathbb{Z}$ satisfy $0\le p\le P$. If executing $i$ from label $l$ with register
vector determined by $st$, no prior answer, and memory configuration $(M,p)$ can reach a
configuration $c$ in finitely many steps, with $G(x,B,P)$ holding at every time before
reaching $c$, then the machine started at label $l$ with register vector determined by
$st$, no prior answer, and memory configuration $(M,p)$ can also reach $c$ in finitely
many steps, with $G(x,B,P)$ holding at every time before reaching $c$.
\end{lemma}

\begin{lemma}[Deterministic Step Preserves Halting]
\lean{Complexity.LogProg.ExpBFS.stepH_of}
\label{Complexity.LogProg.ExpBFS.stepH_of}
\leanok
\uses{Complexity.LogProg.ExpBFS.A, Complexity.LogProg.ExpBFS.G, Complexity.LogProg.ExpBFS.Lb, Complexity.LogProg.ExpBFS.St, Complexity.LogProg.ExpBFS.St.Bd, Complexity.LogProg.ExpBFS.St.toV, Complexity.LogProg.ExpBFS.cf, Complexity.LogProg.ExpBFS.g_of, Complexity.LogProg.Ins, Complexity.LogProg.MHalt, Complexity.LogProg.MHalt.step, Complexity.LogProg.astep, Complexity.LogProg.mstep}
\difficulty{3}
Let $K_0,B,P\in\mathbb{N}$, let $\mathcal O$ assign a Boolean oracle on finite bit
strings to each of four indices, and let $x$ be a finite bit string. Let $l$ be a
machine label with deterministic instruction $i$, let $st$ be a state whose 17 register
values are all at most $B$, let $M\colon\mathbb{Z}\to\{0,1\}$ be a memory tape, and let
$p\in\mathbb{Z}$ satisfy $0\le p\le P$. If executing $i$ from label $l$ with register
vector determined by $st$, no prior answer, and memory configuration $(M,p)$ yields a
configuration from which the machine eventually halts with answer $r\in\{0,1\}$, with
the guard condition $G(x,B,P)$ holding at every time before halting, then the machine
started at label $l$ with register vector determined by $st$, no prior answer, and
memory configuration $(M,p)$ also eventually halts with answer $r$, with $G(x,B,P)$
holding at every time before halting.
\end{lemma}

\begin{lemma}[Prepending a Right Memory Move]
\lean{Complexity.LogProg.ExpBFS.mright_of}
\label{Complexity.LogProg.ExpBFS.mright_of}
\leanok
\uses{Complexity.LogProg.ExpBFS.A, Complexity.LogProg.ExpBFS.G, Complexity.LogProg.ExpBFS.Lb, Complexity.LogProg.ExpBFS.St, Complexity.LogProg.ExpBFS.St.Bd, Complexity.LogProg.ExpBFS.cf, Complexity.LogProg.ExpBFS.g_of, Complexity.LogProg.MConf, Complexity.LogProg.MReach, Complexity.LogProg.MReach.step, Complexity.LogProg.mstep}
\difficulty{3}
Let $K_0,B,P$ be natural numbers, let
$\mathrm{orc}:\{0,1,2,3\}\to\{0,1\}^*\to\{\mathrm{false},\mathrm{true}\}$ be an oracle,
and let $x$ be a Boolean word. Let $st$ be a machine state whose 17 register values are
all at most $B$, let $M:\mathbb Z\to\{\mathrm{false},\mathrm{true}\}$ be a memory, let
$p\in\mathbb Z$ satisfy $0\le p\le P$, and let $l,l'$ be control states. If the
instruction at $l$ moves the memory head one cell right and continues at $l'$, then
every guarded finite run from the configuration at $l'$ with state $st$, memory $M$, and
head position $p+1$ to a configuration $c$ can be prefixed by that instruction: there is
a guarded finite run from the configuration at $l$ with the same state and memory and
head position $p$ to $c$.
\end{lemma}

\begin{lemma}[Prepending a Left Memory Move]
\lean{Complexity.LogProg.ExpBFS.mleft_of}
\label{Complexity.LogProg.ExpBFS.mleft_of}
\leanok
\uses{Complexity.LogProg.ExpBFS.A, Complexity.LogProg.ExpBFS.G, Complexity.LogProg.ExpBFS.Lb, Complexity.LogProg.ExpBFS.St, Complexity.LogProg.ExpBFS.St.Bd, Complexity.LogProg.ExpBFS.cf, Complexity.LogProg.ExpBFS.g_of, Complexity.LogProg.MConf, Complexity.LogProg.MReach, Complexity.LogProg.MReach.step, Complexity.LogProg.mstep}
\difficulty{3}
Let $K_0,B,P$ be natural numbers, let
$\mathrm{orc}:\{0,1,2,3\}\to\{0,1\}^*\to\{\mathrm{false},\mathrm{true}\}$ be an oracle,
and let $x$ be a Boolean word. Let $st$ be a machine state whose 17 register values are
all at most $B$, let $M:\mathbb Z\to\{\mathrm{false},\mathrm{true}\}$ be a memory, let
$p\in\mathbb Z$ satisfy $0\le p\le P$, and let $l,l'$ be control states. If the
instruction at $l$ moves the memory head one cell left and continues at $l'$, then every
guarded finite run from the configuration at $l'$ with state $st$, memory $M$, and head
position $p-1$ to a configuration $c$ can be prefixed by that instruction: there is a
guarded finite run from the configuration at $l$ with the same state and memory and head
position $p$ to $c$.
\end{lemma}

\begin{lemma}[Preservation across setting a memory bit]
\lean{Complexity.LogProg.ExpBFS.mset_of}
\label{Complexity.LogProg.ExpBFS.mset_of}
\leanok
\uses{Complexity.LogProg.ExpBFS.A, Complexity.LogProg.ExpBFS.G, Complexity.LogProg.ExpBFS.Lb, Complexity.LogProg.ExpBFS.St, Complexity.LogProg.ExpBFS.St.Bd, Complexity.LogProg.ExpBFS.cf, Complexity.LogProg.ExpBFS.g_of, Complexity.LogProg.MConf, Complexity.LogProg.MReach, Complexity.LogProg.MReach.step, Complexity.LogProg.mstep}
\difficulty{3}
Let $K_0,B,P$ be natural numbers, let
$\mathrm{orc}:\{0,\ldots,3\}\times\{0,1\}^*\to\{\mathrm{false},\mathrm{true}\}$ be an
oracle, and let $x$ be a Boolean word. Suppose the machine instruction at state $l$ sets
the memory bit at the head to true and continues at $l'$. Let $st$ be a register state
with every one of its 17 counters at most $B$, let
$M:\mathbb{Z}\to\{\mathrm{false},\mathrm{true}\}$ be the memory contents, and suppose
the head position $p$ satisfies $0\le p\le P$. If the machine can reach a configuration
$c$ from state $l'$ with state $st$, memory updated to true at $p$, and head at $p$,
with every configuration strictly before the endpoint satisfying the guard, then it can
reach the same $c$ from state $l$ with state $st$, original memory $M$, and head at $p$,
again satisfying the guard at every configuration strictly before the endpoint. The
guard requires that all register values are at most $B$, the head position lies in
$[0,P]$, and the current instruction meets its operand conditions.
\end{lemma}

\begin{lemma}[Preservation Across a Memory Read]
\lean{Complexity.LogProg.ExpBFS.mread_of}
\label{Complexity.LogProg.ExpBFS.mread_of}
\leanok
\uses{Complexity.LogProg.ExpBFS.A, Complexity.LogProg.ExpBFS.G, Complexity.LogProg.ExpBFS.Lb, Complexity.LogProg.ExpBFS.St, Complexity.LogProg.ExpBFS.St.Bd, Complexity.LogProg.ExpBFS.cf, Complexity.LogProg.ExpBFS.g_of, Complexity.LogProg.MConf, Complexity.LogProg.MReach, Complexity.LogProg.MReach.step, Complexity.LogProg.mstep}
\difficulty{3}
Let $K_0,B,P$ be natural numbers, let
$\mathrm{orc}:\{0,1,2,3\}\to\{0,1\}^*\to\{\mathrm{false},\mathrm{true}\}$ be an oracle,
and let $x$ be a Boolean word. Let $st$ be a machine state whose 17 register values are
all at most $B$, let $M:\mathbb Z\to\{\mathrm{false},\mathrm{true}\}$ be a memory, let
$p\in\mathbb Z$ satisfy $0\le p\le P$, and let $l,l_1,l_0$ be control states. If the
instruction at $l$ reads the memory bit at position $p$, continuing at $l_1$ when $M(p)$
is true and at $l_0$ when it is false, then every guarded finite run from the selected
successor configuration—with state $st$, memory $M$, and head position $p$—to a
configuration $c$ can be prefixed by this read: there is a guarded finite run from the
configuration at $l$ with the same state, memory, and head position to $c$.
\end{lemma}

\begin{lemma}[Reachability under equal register states]
\lean{Complexity.LogProg.ExpBFS.cont}
\label{Complexity.LogProg.ExpBFS.cont}
\leanok
\uses{Complexity.LogProg.ExpBFS.A, Complexity.LogProg.ExpBFS.G, Complexity.LogProg.ExpBFS.Lb, Complexity.LogProg.ExpBFS.St, Complexity.LogProg.ExpBFS.cf, Complexity.LogProg.MConf, Complexity.LogProg.MReach}
\difficulty{1}
Let $K_0,B,P$ be natural numbers, let $\mathsf{orc}$ assign a Boolean to each of four
query types and each finite bit string, and let $x$ be a finite bit string. For machine
states $st$ and $st'$, if $st=st'$ and there is a finite run satisfying the guard from
the configuration with control state $l$, register vector encoded by $st'$, memory
$M:\mathbb Z\to\{\mathrm{false},\mathrm{true}\}$, and head position $p\in\mathbb Z$ to a
configuration $c$, then there is such a run from the configuration with the same control
state, memory, and head position and register vector encoded by $st$, to $c$.
\end{lemma}

\begin{lemma}[Lifting a Register Run to the Memory Machine]
\lean{Complexity.LogProg.ExpBFS.of_areach}
\label{Complexity.LogProg.ExpBFS.of_areach}
\leanok
\uses{Complexity.LogProg.AConf, Complexity.LogProg.AReach, Complexity.LogProg.ExpBFS.A, Complexity.LogProg.ExpBFS.Ad, Complexity.LogProg.ExpBFS.G, Complexity.LogProg.ExpBFS.Lb, Complexity.LogProg.GuardLe, Complexity.LogProg.MReach, Complexity.LogProg.MReach.mono, Complexity.LogProg.MReach.of_areach}
\difficulty{3}
Let $K_0,B,P$ be natural numbers, let $\mathrm{orc}$ be a Boolean oracle on four indices
and Boolean strings, and let $x$ be a Boolean string. Let $a,a'$ be register
configurations, let $M:\mathbb{Z}\to\{\mathrm{false},\mathrm{true}\}$ be memory, and let
$p\in\mathbb{Z}$. Suppose $0\le p\le P$, the register machine reaches $a'$ from $a$ with
the guard $\mathrm{GuardLe}(\mathrm{Ad}(K_0),x,B)$ holding at every configuration
strictly before the endpoint, and the control state of $a'$ is nonhalting. Then the
memory machine reaches $(a',(M,p))$ from $(a,(M,p))$, with $G(x,B,P)$ holding at every
configuration strictly before the endpoint.
\end{lemma}

\begin{lemma}[Reachability across a cleared register]
\lean{Complexity.LogProg.ExpBFS.clr_step}
\label{Complexity.LogProg.ExpBFS.clr_step}
\leanok
\uses{Complexity.LogProg.ExpBFS.A, Complexity.LogProg.ExpBFS.G, Complexity.LogProg.ExpBFS.Lb, Complexity.LogProg.ExpBFS.St, Complexity.LogProg.ExpBFS.St.Bd, Complexity.LogProg.ExpBFS.St.toV, Complexity.LogProg.ExpBFS.cf, Complexity.LogProg.ExpBFS.step_of, Complexity.LogProg.MConf, Complexity.LogProg.MReach}
\difficulty{1}
Let $K_0,B,P$ be natural numbers, let $\mathsf{orc}$ assign a Boolean to each of four
query types and each finite bit string, and let $x$ be a finite bit string. Let $st$ be
a machine state whose 17 register values are all at most $B$, let $M:\mathbb
Z\to\{\mathrm{false},\mathrm{true}\}$ be a memory function, and let $p\in\mathbb Z$
satisfy $0\le p\le P$. If the machine's instruction at state $l$ clears register $r$ and
moves to state $l'$, and there is a finite guarded run from state $l'$ with register
vector obtained from $st$ by setting register $r$ to zero, memory $M$, and memory head
at $p$, to a configuration $c$, then there is such a guarded run from state $l$ with
register vector $st$, the same memory and head position, to $c$.
\end{lemma}

\begin{lemma}[Reachability Before an Increment]
\lean{Complexity.LogProg.ExpBFS.inc_step}
\label{Complexity.LogProg.ExpBFS.inc_step}
\leanok
\uses{Complexity.LogProg.ExpBFS.A, Complexity.LogProg.ExpBFS.G, Complexity.LogProg.ExpBFS.Lb, Complexity.LogProg.ExpBFS.St, Complexity.LogProg.ExpBFS.St.Bd, Complexity.LogProg.ExpBFS.St.toV, Complexity.LogProg.ExpBFS.cf, Complexity.LogProg.ExpBFS.step_of, Complexity.LogProg.MConf, Complexity.LogProg.MReach}
\difficulty{1}
Let $K_0,B,P$ be natural numbers, let $\mathrm{orc}$ be an oracle on four indices and
Boolean strings, and let $x$ be a Boolean string. Suppose an instruction at label $l$
increments register $r$ and continues at $l'$. Let $S$ be a machine state, $M:\mathbb
Z\to\{\mathrm{false},\mathrm{true}\}$ a memory, $p\in\mathbb Z$, and $c$ a
configuration. If every coordinate of $S$ is at most $B$, $0\le p\le P$, and there is a
run from label $l'$ with state $S$ updated by increasing register $r$ by one, empty
result, memory $M$, and head at $p$ to $c$, satisfying $G(x,B,P)$ at every earlier step,
then there is such a run from label $l$ with the original state $S$, memory $M$, and
head at $p$ to $c$, again satisfying $G(x,B,P)$ at every earlier step.
\end{lemma}

\begin{lemma}[Distinct copy registers]
\lean{Complexity.LogProg.ExpBFS.cp_ne}
\label{Complexity.LogProg.ExpBFS.cp_ne}
\leanok
\uses{Complexity.LogProg.ExpBFS.CpS, Complexity.LogProg.ExpBFS.Site.idx, Complexity.LogProg.ExpBFS.cpDst, Complexity.LogProg.ExpBFS.cpSrc, Complexity.LogProg.ExpBFS.rZ, Complexity.LogProg.ExpBFS.rad}
\difficulty{4}
For every copy site $k$, the source and destination register indices are distinct, and
neither is the scratch register index $12$.
\end{lemma}

\begin{theorem}[Copy routine run]
\lean{Complexity.LogProg.ExpBFS.cp_run}
\label{Complexity.LogProg.ExpBFS.cp_run}
\leanok
\uses{Complexity.LogProg.ExpBFS.A, Complexity.LogProg.ExpBFS.Ad, Complexity.LogProg.ExpBFS.CpS, Complexity.LogProg.ExpBFS.G, Complexity.LogProg.ExpBFS.St, Complexity.LogProg.ExpBFS.St.Bd, Complexity.LogProg.ExpBFS.St.Bd.toV, Complexity.LogProg.ExpBFS.St.toV, Complexity.LogProg.ExpBFS.cf, Complexity.LogProg.ExpBFS.cpDst, Complexity.LogProg.ExpBFS.cpExit, Complexity.LogProg.ExpBFS.cpSrc, Complexity.LogProg.ExpBFS.cp_ne, Complexity.LogProg.ExpBFS.of_areach, Complexity.LogProg.MReach, Complexity.LogProg.copy_run}
\difficulty{1}
Let $K_0,B,P$ be natural numbers, let $k$ be a copy site, and let $S$ be a machine
state, $M\colon\mathbb{Z}\to\{\mathrm{false},\mathrm{true}\}$ a memory configuration,
and $p\in\mathbb{Z}$. If $0\le p\le P$, every register value of $S$ is at most $B$, its
scratch register $Z$ is zero, and the sum of the source and destination register values
for $k$ is at most $B$, then the machine has a finite guarded run from the first phase
of the copy routine, initialized with $S$, memory $M$, and head position $p$, to the
exit control state for $k$. At the endpoint, the destination register contains the sum
of its initial value and the source register's initial value; all other registers, the
memory, and the head position are unchanged.
\end{theorem}

\begin{theorem}[Shift Run]
\lean{Complexity.LogProg.ExpBFS.sh_run}
\label{Complexity.LogProg.ExpBFS.sh_run}
\leanok
\uses{Complexity.LogProg.ExpBFS.A, Complexity.LogProg.ExpBFS.Ad, Complexity.LogProg.ExpBFS.G, Complexity.LogProg.ExpBFS.Ph3, Complexity.LogProg.ExpBFS.St, Complexity.LogProg.ExpBFS.St.Bd, Complexity.LogProg.ExpBFS.St.Bd.toV, Complexity.LogProg.ExpBFS.St.toV, Complexity.LogProg.ExpBFS.cf, Complexity.LogProg.ExpBFS.of_areach, Complexity.LogProg.ExpBFS.rC, Complexity.LogProg.ExpBFS.rK, Complexity.LogProg.ExpBFS.rZ, Complexity.LogProg.MReach, Complexity.LogProg.shl_run}
\difficulty{3}
Let $K_0,B,P$ be natural numbers, let $\mathcal O$ assign a Boolean to each of four
query indices and each finite bit string, let $x$ be a finite bit string, let $q$ be one
of the three search phases, let $S$ be a machine state, let $M:\mathbb
Z\to\{\mathrm{false},\mathrm{true}\}$ be a memory function, and let $p\in\mathbb Z$.
Suppose $0\le p\le P$, every register of $S$ is at most $B$, $S.Z=0$, and
$S.K\,2^{S.C}\le B$. Then the machine has a finite run from shift phase $0$ for phase
$q$ to the first copy phase for that same phase, with every configuration satisfying
$G(x,B,P)$. At the endpoint the shift counter is zero, $K=S.K\,2^{S.C}$ (using the
initial values in $S$), and all other registers, the memory function, and the
memory-head position are unchanged.
\end{theorem}

\begin{theorem}[Counter transfer in the breadth-first search machine]
\lean{Complexity.LogProg.ExpBFS.mv_run}
\label{Complexity.LogProg.ExpBFS.mv_run}
\leanok
\uses{Complexity.LogProg.ExpBFS.A, Complexity.LogProg.ExpBFS.Ad, Complexity.LogProg.ExpBFS.G, Complexity.LogProg.ExpBFS.St, Complexity.LogProg.ExpBFS.St.Bd, Complexity.LogProg.ExpBFS.St.Bd.toV, Complexity.LogProg.ExpBFS.St.toV, Complexity.LogProg.ExpBFS.cf, Complexity.LogProg.ExpBFS.of_areach, Complexity.LogProg.ExpBFS.rc, Complexity.LogProg.ExpBFS.rc1, Complexity.LogProg.MReach, Complexity.LogProg.move_run}
\difficulty{3}
Let $K_0,B,P$ be natural numbers, let $x$ be a finite Boolean word, and let
$\mathrm{orc}$ assign a Boolean answer to each of four indices and each finite Boolean
word. Let $st$ be a machine state whose 17 natural-number counters are all at most $B$,
let $M:\mathbb Z\to\{\mathrm{false},\mathrm{true}\}$ be a memory map, and let
$p\in\mathbb Z$ satisfy $0\le p\le P$ and $st.c_1+st.c\le B$. In the machine defined by
the transition rules of the exponential breadth-first search procedure, the
configuration starting in the counter-transfer phase with state $st$, memory $M$, and
memory-head position $p$ reaches the copy phase for the radius and next-radius counters
with the same memory and head position, with counter $c$ reset to $0$ and counter $c_1$
set to $st.c_1+st.c$. Every configuration strictly before that endpoint has all register
values at most $B$, has memory-head position between $0$ and $P$, and satisfies the
required operand conditions for its active instruction: equality tests compare distinct
registers, input-word comparisons occur only on valid encodings, and calls use distinct
argument registers.
\end{theorem}

\begin{lemma}[Reachability under equal endpoint data]
\lean{Complexity.LogProg.ExpBFS.reach_cf}
\label{Complexity.LogProg.ExpBFS.reach_cf}
\leanok
\uses{Complexity.LogProg.ExpBFS.A, Complexity.LogProg.ExpBFS.G, Complexity.LogProg.ExpBFS.Lb, Complexity.LogProg.ExpBFS.St, Complexity.LogProg.ExpBFS.cf, Complexity.LogProg.MConf, Complexity.LogProg.MReach}
\difficulty{1}
Fix a natural parameter $K_0$, an oracle, a Boolean string $x$, and natural numbers
$B,P$. If a configuration $a$ reaches the configuration with label $l$, register state
$st$, memory $M$, and head position $p$, with the guard holding at every time strictly
before arrival, then whenever $st=st'$ and $p=p'$, $a$ also reaches the configuration
with label $l$, register state $st'$, memory $M$, and head position $p'$, with the same
guard condition.
\end{lemma}

\begin{lemma}[Register Index at Each Site]
\lean{Complexity.LogProg.ExpBFS.idx_cases}
\label{Complexity.LogProg.ExpBFS.idx_cases}
\leanok
\uses{Complexity.LogProg.ExpBFS.Site, Complexity.LogProg.ExpBFS.Site.idx, Complexity.LogProg.ExpBFS.rU, Complexity.LogProg.ExpBFS.rV}
\difficulty{3}
For every search site, its associated register index is either the index of $U$ or the
index of $V$.
\end{lemma}

\begin{theorem}[Bounded outward walk]
\lean{Complexity.LogProg.ExpBFS.wo_run}
\label{Complexity.LogProg.ExpBFS.wo_run}
\leanok
\uses{Complexity.LogProg.ExpBFS.A, Complexity.LogProg.ExpBFS.G, Complexity.LogProg.ExpBFS.Site, Complexity.LogProg.ExpBFS.St, Complexity.LogProg.ExpBFS.St.Bd, Complexity.LogProg.ExpBFS.St.toV_T, Complexity.LogProg.ExpBFS.cf, Complexity.LogProg.ExpBFS.mright_of, Complexity.LogProg.ExpBFS.reach_cf, Complexity.LogProg.ExpBFS.step_of, Complexity.LogProg.MReach, Complexity.LogProg.MReach.refl, Complexity.LogProg.astep}
\difficulty{6}
For the designated machine with parameter $K_0$, oracle
$\mathrm{orc}:\{0,\ldots,3\}\times\{0,1\}^*\to\{0,1\}$, input $x$, and natural bounds
$B,P$, let $st$ be a memory access site, $k$ a natural number, $S$ a machine state,
$M:\mathbb Z\to\{0,1\}$ a memory marking, and $p\in\mathbb Z$. If $S.T=k$, every
register of $S$ is at most $B$, $S.T_2+k\le B$, $0\le p$, and $p+k\le P$, then the
machine has a finite run from the outward-walk configuration at $st$, with state $S$,
marking $M$, and head position $p$, to the memory-operation configuration at $st$, with
the same marking, head position $p+k$, and state obtained by setting $T$ to zero and
$T_2$ to $S.T_2+k$, leaving every other state component unchanged. Every configuration
strictly before the endpoint satisfies the register-instruction precondition, has all
registers at most $B$, and has memory head between $0$ and $P$.
\end{theorem}

\begin{theorem}[Bounded return walk]
\lean{Complexity.LogProg.ExpBFS.wb_run}
\label{Complexity.LogProg.ExpBFS.wb_run}
\leanok
\uses{Complexity.LogProg.ExpBFS.A, Complexity.LogProg.ExpBFS.G, Complexity.LogProg.ExpBFS.Site, Complexity.LogProg.ExpBFS.St, Complexity.LogProg.ExpBFS.St.Bd, Complexity.LogProg.ExpBFS.St.toV_T2, Complexity.LogProg.ExpBFS.acExit, Complexity.LogProg.ExpBFS.cf, Complexity.LogProg.ExpBFS.mleft_of, Complexity.LogProg.ExpBFS.reach_cf, Complexity.LogProg.ExpBFS.step_of, Complexity.LogProg.MReach, Complexity.LogProg.MReach.refl, Complexity.LogProg.astep}
\difficulty{5}
For the designated machine with parameter $K_0$, oracle
$\mathrm{orc}:\{0,\ldots,3\}\times\{0,1\}^*\to\{0,1\}$, input $x$, and natural bounds
$B,P$, let $st$ be a memory access site, $b$ a Boolean, $k$ a natural number, $S$ a
machine state, $M:\mathbb Z\to\{0,1\}$ a memory marking, and $p\in\mathbb Z$. If
$S.T_2=k$, every register of $S$ is at most $B$, $0\le p-k$, and $p\le P$, then the
machine has a finite run from the return-walk configuration at $st$ carrying bit $b$,
with state $S$, marking $M$, and head position $p$, to the configuration at the
corresponding exit with state $S$ except that $T_2$ is set to zero, the same marking,
and head position $p-k$. At every configuration strictly before the endpoint, the
register-instruction precondition holds, every register is at most $B$, and the memory
head lies between $0$ and $P$.
\end{theorem}

\begin{theorem}[Memory access simulation]
\lean{Complexity.LogProg.ExpBFS.ac_run}
\label{Complexity.LogProg.ExpBFS.ac_run}
\leanok
\uses{Complexity.LogProg.ExpBFS.A, Complexity.LogProg.ExpBFS.G, Complexity.LogProg.ExpBFS.Site, Complexity.LogProg.ExpBFS.Site.idx, Complexity.LogProg.ExpBFS.Site.set, Complexity.LogProg.ExpBFS.St, Complexity.LogProg.ExpBFS.St.Bd, Complexity.LogProg.ExpBFS.St.toV, Complexity.LogProg.ExpBFS.St.toV_T, Complexity.LogProg.ExpBFS.acExit, Complexity.LogProg.ExpBFS.acI, Complexity.LogProg.ExpBFS.cf, Complexity.LogProg.ExpBFS.cpDst, Complexity.LogProg.ExpBFS.cpExit, Complexity.LogProg.ExpBFS.cpSrc, Complexity.LogProg.ExpBFS.cp_run, Complexity.LogProg.ExpBFS.idx_cases, Complexity.LogProg.ExpBFS.mread_of, Complexity.LogProg.ExpBFS.mset_of, Complexity.LogProg.ExpBFS.reach_cf, Complexity.LogProg.ExpBFS.wb_run, Complexity.LogProg.ExpBFS.wo_run, Complexity.LogProg.MReach, Complexity.LogProg.MReach.trans}
\difficulty{6}
Let $K_0,B,P$ be natural numbers, let $\mathcal O$ assign a Boolean answer to each of
four query indices and each finite Boolean word, and let $x$ be a finite Boolean word.
Choose a site $st$, a machine state $S$, and a memory $M\colon\mathbb
Z\to\{\mathrm{false},\mathrm{true}\}$. Assume every counter of $S$ is at most $B$, its
outward-walk, return-walk, and scratch counters are zero, and the address $a$ given by
the memory-region base plus the counter selected by $st$ is at most both $B$ and $P$.
Starting in the access routine for $st$ with state $S$, memory $M$, and memory head at
$0$, the machine reaches the exit specified by $st$ and the bit $st.set\lor M(a)$, with
state again $S$, head again at $0$, and memory changed only by setting cell $a$ to true
when $st$ is a marking site. At every earlier configuration, all counters are at most
$B$, the memory head lies between $0$ and $P$, register comparisons use distinct
registers, input comparisons are made only on inputs with the required valid plain or
paired form, and oracle calls use lists of distinct registers.
\end{theorem}

\begin{definition}[Binary Logarithmic Width]
\lean{Complexity.LogProg.ExpBFS.hv}
\label{Complexity.LogProg.ExpBFS.hv}
\leanok
For a natural number $k$, define $h(k)=0$ if $k=0$, and otherwise define
$h(k)=\lfloor\log_2 k\rfloor+1$.
\end{definition}

\begin{lemma}[Halving the Logarithmic Width]
\lean{Complexity.LogProg.ExpBFS.hv_div}
\label{Complexity.LogProg.ExpBFS.hv_div}
\leanok
\uses{Complexity.LogProg.ExpBFS.hv}
\difficulty{4}
For every positive natural number $k$, the function $h(k)$ defined to be $0$ at $k=0$
and $\lfloor\log_2 k\rfloor+1$ otherwise satisfies $h(k)=h(\lfloor k/2\rfloor)+1$.
\end{lemma}

\begin{definition}[Logarithmic Bound]
\lean{Complexity.LogProg.ExpBFS.lgv}
\label{Complexity.LogProg.ExpBFS.lgv}
\leanok
\uses{Complexity.LogProg.ExpBFS.hv}
For a natural number $n$, define $g(n)=1+h(n)$, where $h(0)=0$ and $h(n)=\lfloor\log_2
n\rfloor+1$ for $n>0$.
\end{definition}

\begin{lemma}[Comparison of logarithmic bounds]
\lean{Complexity.LogProg.ExpBFS.logSpace_le_lgv}
\label{Complexity.LogProg.ExpBFS.logSpace_le_lgv}
\leanok
\uses{Complexity.LogProg.ExpBFS.hv, Complexity.LogProg.ExpBFS.lgv, Complexity.logSpace}
\difficulty{3}
For every $n\in\mathbb N$, the logarithmic space bound $\lfloor\log_2 n\rfloor+1$ is at
most $1+h(n)$, where $h(0)=0$ and $h(n)=\lfloor\log_2 n\rfloor+1$ for $n>0$.
\end{lemma}

\begin{lemma}[Upper bound for a logarithmic quantity]
\lean{Complexity.LogProg.ExpBFS.lgv_le}
\label{Complexity.LogProg.ExpBFS.lgv_le}
\leanok
\uses{Complexity.LogProg.ExpBFS.hv, Complexity.LogProg.ExpBFS.lgv, Complexity.logSpace}
\difficulty{3}
For every natural number $n$, the quantity equal to $1$ when $n=0$ and to
$2+\lfloor\log_2 n\rfloor$ when $n>0$ is at most $2+\lfloor\log_2 n\rfloor$.
\end{lemma}

\begin{lemma}[Linear upper bound for the adjusted logarithm]
\lean{Complexity.LogProg.ExpBFS.lgv_le_succ}
\label{Complexity.LogProg.ExpBFS.lgv_le_succ}
\leanok
\uses{Complexity.LogProg.ExpBFS.hv, Complexity.LogProg.ExpBFS.lgv}
\difficulty{3}
For every natural number $n$, the adjusted logarithmic quantity $lgv(n)$ is at most
$n+1$.
\end{lemma}

\begin{theorem}[Logarithm setup reaches its initialized state]
\lean{Complexity.LogProg.ExpBFS.lg_run}
\label{Complexity.LogProg.ExpBFS.lg_run}
\leanok
\uses{Complexity.LogProg.ExpBFS.A, Complexity.LogProg.ExpBFS.G, Complexity.LogProg.ExpBFS.St, Complexity.LogProg.ExpBFS.St.Bd, Complexity.LogProg.ExpBFS.cf, Complexity.LogProg.ExpBFS.hv, Complexity.LogProg.ExpBFS.hv_div, Complexity.LogProg.ExpBFS.reach_cf, Complexity.LogProg.ExpBFS.step_of, Complexity.LogProg.MReach, Complexity.LogProg.MReach.refl, Complexity.LogProg.astep}
\difficulty{6}
Fix a natural number $K_0$, an oracle assigning a Boolean answer to each of four query
types and each Boolean string, and a Boolean string $x$. For every natural number $k$,
state $S$, and Boolean memory $M:\mathbb{Z}\to\{\mathrm{false},\mathrm{true}\}$, if
$S.N=k$, every register of $S$ is at most $B$, and $S.Lg+hv(k)\le B$, then the program
$A(K_0)$ reaches its initialized search state with input-length counter $0$ and
logarithm counter $S.Lg+hv(k)$. Every configuration strictly before the endpoint
satisfies the guard with register bound $B$ and memory-head position between $0$ and
$P$, inclusive; the memory remains $M$ and the head is at $0$ at the endpoint.
\end{theorem}

\begin{theorem}[Initial Setup Run]
\lean{Complexity.LogProg.ExpBFS.setup_run}
\label{Complexity.LogProg.ExpBFS.setup_run}
\leanok
\uses{Complexity.LogProg.ExpBFS.A, Complexity.LogProg.ExpBFS.Ad, Complexity.LogProg.ExpBFS.G, Complexity.LogProg.ExpBFS.St.Bd, Complexity.LogProg.ExpBFS.St.Bd.toV, Complexity.LogProg.ExpBFS.St.toV, Complexity.LogProg.ExpBFS.St.zero, Complexity.LogProg.ExpBFS.cf, Complexity.LogProg.ExpBFS.hv, Complexity.LogProg.ExpBFS.lg_run, Complexity.LogProg.ExpBFS.lgv, Complexity.LogProg.ExpBFS.of_areach, Complexity.LogProg.ExpBFS.rN, Complexity.LogProg.ExpBFS.reach_cf, Complexity.LogProg.ExpBFS.step_of, Complexity.LogProg.MReach, Complexity.LogProg.MReach.trans, Complexity.LogProg.astep, Complexity.LogProg.inputLength_loop, Turing.pairEncode}
\difficulty{5}
Let $K_0$ be a natural number, let $\mathcal O$ assign a Boolean to each of four query
indices and each finite bit string, let $x$ be a finite bit string, and let $M:\mathbb
Z\to\{\mathrm{false},\mathrm{true}\}$ be a memory function. Suppose that for every $k\le
|x|$, the first oracle answers true on the encoding of $x$ paired with the binary
representation of $k$ exactly when $k<|x|$, and suppose $|x|+2\le B$. Then the machine
has a finite run from the initial scan configuration with memory $M$ and head position
$0$ to the first radius-copy phase, with every configuration satisfying the register and
memory-head bounds in $G(x,B,P)$. At the endpoint the memory and head position are still
$M$ and $0$, respectively, and all registers are zero except the logarithm-bound
register, which equals $1+h(|x|)$, where $h(n)=0$ for $n=0$ and $h(n)=\lfloor\log_2
n\rfloor+1$ for $n>0$.
\end{theorem}

\begin{theorem}[Repeated register addition run]
\lean{Complexity.LogProg.ExpBFS.mw_run}
\label{Complexity.LogProg.ExpBFS.mw_run}
\leanok
\uses{Complexity.LogProg.ExpBFS.A, Complexity.LogProg.ExpBFS.G, Complexity.LogProg.ExpBFS.Ph3, Complexity.LogProg.ExpBFS.St, Complexity.LogProg.ExpBFS.St.Bd, Complexity.LogProg.ExpBFS.cf, Complexity.LogProg.ExpBFS.cpDst, Complexity.LogProg.ExpBFS.cpExit, Complexity.LogProg.ExpBFS.cpSrc, Complexity.LogProg.ExpBFS.cp_run, Complexity.LogProg.ExpBFS.reach_cf, Complexity.LogProg.MReach, Complexity.LogProg.MReach.trans}
\difficulty{5}
Let $K_0,B,P$ be natural numbers, let $\mathrm{orc}$ be a Boolean-valued oracle on four
indices and finite Boolean strings, and let $x$ be a Boolean string. Choose a phase $q$
among the first search, second search, and final scan, natural numbers $j,i$ with
$i<K_0$ and $K_0-i=j$, a machine state $S$, and memory contents
$M:\mathbb{Z}\to\{\mathrm{false},\mathrm{true}\}$. Suppose the width register of $S$ is
$i$ times its width factor, its scratch register is zero, every register is at most $B$,
and $K_0$ times its width factor is at most $B$. From the configuration at copy-loop
stage $0$ for index $i$ and phase $q$, with registers $S$, memory $M$, and memory head
at $0$, the machine reaches the instruction to clear the width register with that
register equal to $K_0$ times the original width factor, all other registers unchanged,
memory still $M$, and memory head still at $0$; at every step strictly before the
endpoint, the register side conditions hold, every register is at most $B$, and the
memory head lies between $0$ and $P$, inclusive.
\end{theorem}

\begin{definition}[Width-Adjusted Machine State]
\lean{Complexity.LogProg.ExpBFS.wkSt}
\label{Complexity.LogProg.ExpBFS.wkSt}
\leanok
\uses{Complexity.LogProg.ExpBFS.Ph3, Complexity.LogProg.ExpBFS.St}
Given a phase $q$, a state $S$, and a natural number $r$, define a new state by
replacing its code width $w$ with $K_0(r+S.Lg+1)$, its code bound $K$ with
$2^{K_0(r+S.Lg+1)}$, and its memory base $Bs$ with $2^{K_0(r+S.Lg+1)}$ when $q$ is the
second-search phase and twice that value otherwise. All other components remain those of
$S$.
\end{definition}

\begin{theorem}[Bounded width setup run]
\lean{Complexity.LogProg.ExpBFS.wk_run}
\label{Complexity.LogProg.ExpBFS.wk_run}
\leanok
\uses{Complexity.LogProg.ExpBFS.A, Complexity.LogProg.ExpBFS.G, Complexity.LogProg.ExpBFS.Ph3, Complexity.LogProg.ExpBFS.St, Complexity.LogProg.ExpBFS.St.Bd, Complexity.LogProg.ExpBFS.St.toV, Complexity.LogProg.ExpBFS.cf, Complexity.LogProg.ExpBFS.cont, Complexity.LogProg.ExpBFS.cpDst, Complexity.LogProg.ExpBFS.cpExit, Complexity.LogProg.ExpBFS.cpSrc, Complexity.LogProg.ExpBFS.cp_run, Complexity.LogProg.ExpBFS.mw_run, Complexity.LogProg.ExpBFS.rad, Complexity.LogProg.ExpBFS.reach_cf, Complexity.LogProg.ExpBFS.sh_run, Complexity.LogProg.ExpBFS.step_of, Complexity.LogProg.ExpBFS.wkExit, Complexity.LogProg.ExpBFS.wkSt, Complexity.LogProg.MReach, Complexity.LogProg.MReach.refl, Complexity.LogProg.MReach.trans, Complexity.LogProg.astep}
\difficulty{5}
For a natural parameter $K_0$, the designated machine, oracle
$\mathrm{orc}:\{0,\ldots,3\}\times\{0,1\}^*\to\{0,1\}$, input $x$, and natural bounds
$B,P$, let $q$ be one of the first-search, second-search, or final-scan phases, let $S$
be a machine state, and let $M:\mathbb Z\to\{0,1\}$ be a memory marking. Suppose
$X=w=K=B_s=C=Z=0$ in $S$, $K_0>0$, every register of $S$ is at most $B$, and $3\cdot
2^{K_0(r+S.L_g+1)}\le B$, where $r$ is the radius register selected by phase $q$. Then
the machine has a finite run from the copy-loop configuration for the radius of phase
$q$, with state $S$, marking $M$, and memory head at zero, to the phase exit with the
same marking and head position zero, and with state obtained by setting
$w=K_0(r+S.L_g+1)$, $K=2^{K_0(r+S.L_g+1)}$, and $B_s=(1\text{ if }q\text{ is the
second-search phase, and }2\text{ otherwise})\cdot 2^{K_0(r+S.L_g+1)}$, leaving every
other state component unchanged. Every configuration strictly before the endpoint
satisfies the register-instruction precondition, has all registers at most $B$, and has
memory head between $0$ and $P$.
\end{theorem}

\begin{theorem}[Register reset run]
\lean{Complexity.LogProg.ExpBFS.d1_run}
\label{Complexity.LogProg.ExpBFS.d1_run}
\leanok
\uses{Complexity.LogProg.ExpBFS.A, Complexity.LogProg.ExpBFS.G, Complexity.LogProg.ExpBFS.St, Complexity.LogProg.ExpBFS.St.Bd, Complexity.LogProg.ExpBFS.cf, Complexity.LogProg.ExpBFS.clr_step, Complexity.LogProg.ExpBFS.rBs, Complexity.LogProg.ExpBFS.rK, Complexity.LogProg.ExpBFS.rU, Complexity.LogProg.ExpBFS.rV, Complexity.LogProg.ExpBFS.rw, Complexity.LogProg.MReach, Complexity.LogProg.MReach.refl}
\difficulty{4}
Let $K_0,B,P$ be natural numbers, let $S$ be a machine state, and let
$M\colon\mathbb{Z}\to\{\mathrm{false},\mathrm{true}\}$ be a memory configuration. If
every register value of $S$ is at most $B$, then the machine has a finite guarded run
from the first reset control state, initialized with $S$, memory $M$, and head position
$0$, to the move control state. At the endpoint, the registers $U,V,w,K,B_s$ are zero,
and every other register retains its value from $S$; the memory and head position are
unchanged.
\end{theorem}

\begin{theorem}[Reset Before the Final Scan]
\lean{Complexity.LogProg.ExpBFS.fz_run}
\label{Complexity.LogProg.ExpBFS.fz_run}
\leanok
\uses{Complexity.LogProg.ExpBFS.A, Complexity.LogProg.ExpBFS.G, Complexity.LogProg.ExpBFS.St, Complexity.LogProg.ExpBFS.St.Bd, Complexity.LogProg.ExpBFS.cf, Complexity.LogProg.ExpBFS.clr_step, Complexity.LogProg.ExpBFS.rBs, Complexity.LogProg.ExpBFS.rK, Complexity.LogProg.ExpBFS.rU, Complexity.LogProg.ExpBFS.rV, Complexity.LogProg.ExpBFS.rw, Complexity.LogProg.MReach, Complexity.LogProg.MReach.refl}
\difficulty{4}
Let $K_0,B,P$ be natural numbers, let $x$ be a finite Boolean word, let $\mathrm{orc}$
assign a Boolean answer to each of four indices and each finite Boolean word, and let
$M\colon\mathbb{Z}\to\{\mathrm{false},\mathrm{true}\}$ be a memory marking. For every
machine state $S$ whose 17 register values are each at most $B$, there is a finite run
from the reset stage before the final scan, with registers given by $S$, memory $M$, and
memory head at $0$, to the final-scan preparation stage, with the registers
$U,V,w,K,B_s$ set to zero, all other registers unchanged, memory still $M$, and head
still at $0$. At every time strictly before the endpoint, the register and instruction
safety conditions hold, and the memory head lies between $0$ and $P$, inclusive.
\end{theorem}

\begin{theorem}[Next Iteration Initialization Run]
\lean{Complexity.LogProg.ExpBFS.nz_run}
\label{Complexity.LogProg.ExpBFS.nz_run}
\leanok
\uses{Complexity.LogProg.ExpBFS.A, Complexity.LogProg.ExpBFS.G, Complexity.LogProg.ExpBFS.St, Complexity.LogProg.ExpBFS.St.Bd, Complexity.LogProg.ExpBFS.cf, Complexity.LogProg.ExpBFS.clr_step, Complexity.LogProg.ExpBFS.inc_step, Complexity.LogProg.ExpBFS.rBs, Complexity.LogProg.ExpBFS.rK, Complexity.LogProg.ExpBFS.rU, Complexity.LogProg.ExpBFS.rV, Complexity.LogProg.ExpBFS.rc, Complexity.LogProg.ExpBFS.rc1, Complexity.LogProg.ExpBFS.rs, Complexity.LogProg.ExpBFS.rs1, Complexity.LogProg.ExpBFS.rw, Complexity.LogProg.MReach, Complexity.LogProg.MReach.refl}
\difficulty{5}
Let $K_0,B,P$ be natural numbers, let
$\mathrm{orc}:\{0,1,2,3\}\times\{0,1\}^*\to\{0,1\}$ be an oracle, let $x$ be a finite
Boolean string, let $S$ be a machine state, and let $M:\mathbb{Z}\to\{0,1\}$ be a memory
function. Suppose every register value in $S$ is at most $B$ and $S.s+1\le B$. Starting
in the next-iteration initialization state with registers $S$, memory $M$, and memory
head at $0$, the machine reaches a configuration in the first-search copy phase with
radius $S.s+1$, registers $U,V,w,K,B_s,c,c_1,s_1$ all zero, memory $M$, and head at $0$;
at every earlier configuration, all register values are at most $B$, the memory head
lies between $0$ and $P$, and the current instruction satisfies the machine's
register-comparison side condition.
\end{theorem}
